% =============================================================================
% ANNEXE B — API HTTP
% =============================================================================
\chapter{Catalogue des API HTTP}
\label{annexe:api}

Cette annexe consigne les points d'entrée HTTP exposés par les démons et le contrôleur. Trois ports sont impliqués : 9100 (TLS, binaire, déjà décrit en annexe~\ref{annexe:protocole}), 9200 (HTTP, cluster), 9300 (HTTP, control local).

\section{Port 9200 : API cluster}

Cette API est consommée par le contrôleur Python pour obtenir l'état de chaque nœud. Elle est en HTTP plain car restreinte au VLAN management par le firewall hôte.

\begin{table}[H]
\centering
\small
\caption{Routes du port 9200 (API cluster).}
\begin{tabular}{lp{2.5cm}p{8cm}}
\toprule
\textbf{Route} & \textbf{Méthode} & \textbf{Réponse} \\
\midrule
\texttt{/status}      & GET & JSON \{node\_id, mem\_total\_mib, mem\_free\_mib, pages\_stored, vcpu\_total, vcpu\_free, ceph\_enabled, vms: [\{vmid, avg\_cpu\_pct, throttle\_ratio, remote\_pages\}], gpu: \{backend, vram\_total\_mib, vram\_free\_mib, budgets\}\} \\
\bottomrule
\end{tabular}

\end{table}

\section{Port 9300 : API control local}

Cette API n'est accessible que depuis \texttt{127.0.0.1}. Elle pilote le démon localement et expose les métriques Prometheus.

\begin{table}[H]
\centering
\small
\caption{Routes du port 9300 (API control local).}
\begin{tabular}{lp{2.5cm}p{7.5cm}}
\toprule
\textbf{Route} & \textbf{Méthode} & \textbf{Effet} \\
\midrule
\texttt{/control/status}          & GET  & Version détaillée du status \\
\texttt{/control/metrics}         & GET  & Métriques au format Prometheus \\
\texttt{/control/admit}           & POST & Enregistre un \texttt{VmQuota} (body JSON) \\
\texttt{/control/migrate}         & POST & Déclenche une migration (body : \{vmid, target, type\}) \\
\texttt{/control/evict}           & POST & Force une éviction immédiate (body : \{vmid, page\_id\}) \\
\texttt{/control/vm/\{vmid\}/gpu} & POST & Ajuste le budget VRAM (body : \{vram\_mib\}) \\
\texttt{/control/gpu/status}      & GET  & État détaillé du GPU et budgets \\
\texttt{/control/vm/\{vmid\}/vcpu}& POST & Hotplug/unplug vCPU (body : \{count\}) \\
\bottomrule
\end{tabular}

\end{table}

\section{CLI du contrôleur Python}

Le contrôleur s'invoque via Click :

\begin{lstlisting}[basicstyle=\ttfamily\footnotesize, frame=single]
# Statut instantane
python3 -m controller.main status \
    --stores "192.168.1.2:9100,192.168.1.3:9101"

# Boucle de monitoring
python3 -m controller.main monitor \
    --interval 5 \
    --stores "192.168.1.2:9100,192.168.1.3:9101" \
    --threshold-enable 60 \
    --threshold-migrate 85

# Evaluation de la politique sans execution
python3 -m controller.main policy --dry-run \
    --stores "192.168.1.2:9100,192.168.1.3:9101"

# Niveau de log et format
python3 -m controller.main --log-level INFO --log-format json status
\end{lstlisting}
